\section{The Controller That Could Not Read}
\label{sec:i2c-controller-cannot-read}

This section is a worked example of a specific and unglamorous engineering skill: proving
that a peripheral, and not your code, is at fault. It matters because the alternative —
guessing — is unbounded. You can spend days rewriting a driver that was never the problem.

\subsection{The Symptom}

An I2C0 master was configured with the PAC: bus clock enabled, pins routed, clock divider
set, address written, and a read requested. Writes completed. Reads did not: the transfer
never finished, and the state register reported the controller waiting in the transmit path
with the bus already acknowledged as finished for the write half.

The measurement that matters, taken from \codeinline{I2C0->STATE} after a write-then-read
sequence:

\begin{minted}{text}
I2C0 STATE = 0xd7
    BUSY         = 1
    TRANSMITTER  = 1
    NACKED       = 1
    RXDATAV      = 0     <- no receive data will ever arrive
\end{minted}

\codeinline{RXDATAV = 0} is the whole story: the controller accepted the read command and
never received a byte into its buffer. There is nothing to read out, so any polling loop
written against \codeinline{RXDATAV} waits forever.

\subsection{What Was Ruled Out, and How}

A diagnosis is only worth reading if it says what it eliminated and on what evidence. The
table below is the actual list, in the order it was worked through.

\begin{table}[H]
\centering
\begin{tabular}{p{4.6cm}p{9.2cm}}
\hline
\textbf{Hypothesis} & \textbf{Evidence against} \\
\hline
Wrong address & The same address acknowledged during the write half of the transfer. \\
Wrong pins & A bit-banged master on the same two pins read the device immediately. \\
Missing pull-ups & The pins read high with every GPIO configured as input, no internal
                  pull: the pull-ups are on the board. \\
Missing bus clock & \codeinline{CMU\_CLKEN0.I2C0} read back set; with it cleared the
                    peripheral faults instead of running. \\
Bad clock divider & Two dividers, including the slowest available, produced the same
                    \codeinline{STATE}. \\
The device itself & The device answered a bit-banged transaction on the same bus, at the
                   same address, in the same second. \\
\hline
\end{tabular}
\caption{Hypotheses eliminated before concluding that the on-chip controller was at fault.}
\label{tab:i2c-elimination}
\end{table}

Fourteen measurements across these hypotheses produced the same \codeinline{STATE} every
time. That repetition is the point: a single failing attempt is consistent with a bug, a
dozen identical failures along different paths are consistent with the hardware.

\begin{notebox}
    \textbf{The control experiment is the one that decided it.} The bit-banged master used
    the same pins, the same pull-ups and the same device. It read the correct identifier on
    the first attempt. That single result moved the fault from ``our driver'' to ``this
    controller, on this part, in this configuration'', and it took two minutes to run.
\end{notebox}

\subsection{What To Do About It}

There are three honest responses to a peripheral that does not work, and only one of them is
acceptable in a project that has to ship.

\begin{enumerate}
    \item \textbf{Rewrite the driver.} Cheap to try, unbounded to continue. Only justified
          while the control experiment has not been run.
    \item \textbf{Build a transport that works.} Slower, but fully understood, and it turns
          the bus into something you can debug byte by byte. This is what this project did,
          and it is the subject of the next section.
    \item \textbf{Find the missing configuration.} The correct answer when it exists. It does
          not always: this chapter does not claim the controller is broken in general, only
          that it could not complete a read under every configuration that was tried, and
          that a working alternative was available.
\end{enumerate}

\begin{warningbox}
    \textbf{Write down which one you chose, and why.} A project that silently switches to a
    workaround loses the knowledge that it did. This project keeps the fourteen measurements
    in \filepath{docs/04-reference/I2C0-READ-DEFECT.md}, so that a later engineer — or a
    later silicon revision — can pick up the investigation instead of repeating it.
\end{warningbox}

\subsection{Why This Belongs in a Book}

Because it is the normal experience of driving new hardware, and textbooks rarely show it.
The registers are documented, the sequence is correct, and the thing still does not work.
The skill being taught here is not I2C: it is how to reduce a hardware failure to a
statement you can defend, and then how to keep shipping while it is unresolved.
